\documentclass[12pt]{article}

\usepackage[a4paper,margin=1in]{geometry}
\usepackage{setspace}
\onehalfspacing

\usepackage{graphicx}
\usepackage{amsmath,amssymb,amsfonts,amsthm}
\usepackage{bm}
\usepackage{mathtools}
\usepackage{booktabs}
\usepackage{enumitem}
\usepackage{hyperref}

\setlength{\parindent}{0pt}
\setlength{\parskip}{0.6em}

\setlength{\abovedisplayskip}{10pt}
\setlength{\belowdisplayskip}{10pt}
\setlength{\abovedisplayshortskip}{6pt}
\setlength{\belowdisplayshortskip}{6pt}

\newtheorem{proposition}{Proposition}
\newtheorem{remark}{Remark}

\newcommand{\E}{\mathbb{E}}
\newcommand{\KL}{\operatorname{KL}}
\newcommand{\cov}{\operatorname{cov}}
\newcommand{\Var}{\operatorname{Var}}
\newcommand{\tr}{\operatorname{tr}}
\newcommand{\diag}{\operatorname{diag}}
\newcommand{\vecop}{\operatorname{vec}}

\title{Joint Online Gaussian Process with Linear Mean Model\\
for Kronecker-Separable Spatio-Temporal HiPPO-SVGP}
\author{}
\date{March 2026}

\begin{document}
\maketitle

\section{Joint Observation Model}

We consider a spatio-temporal latent response observed at spatial locations
$s_i\in\mathcal{S}\subset\mathbb{R}^{d_s}$ and time points $t_n\in\mathbb{R}$.
The observation at location $(t_n,s_i)$ is modeled as the sum of a linear mean
component, a nonparametric Gaussian process component, and independent Gaussian noise:
\begin{equation}
y(t_n,s_i)
=
\phi(t_n,s_i)^\top \beta
+
f(t_n,s_i)
+
\epsilon_{n,i},
\qquad
\epsilon_{n,i}\sim\mathcal{N}(0,\sigma^2).
\label{eq:joint-observation-scalar}
\end{equation}
Here $\phi(t_n,s_i)\in\mathbb{R}^{d}$ is a covariate design vector,
$\beta\in\mathbb{R}^{d}$ is the linear mean parameter, and
$f$ is a spatio-temporal Gaussian process.

Let the data in the $n$-th online temporal block be
\begin{equation}
D_n=\{(t,s,y,\phi): (t,s)\in\mathcal{B}_n\},
\end{equation}
and let
\begin{equation}
y_n\in\mathbb{R}^{B_n},
\qquad
\Phi_n\in\mathbb{R}^{B_n\times d}
\end{equation}
denote the stacked response vector and design matrix for this block. The block-level
observation model is
\begin{equation}
y_n = \Phi_n \beta + f_n + \epsilon_n,
\qquad
\epsilon_n\sim\mathcal{N}(0,\sigma^2 I).
\label{eq:block-observation}
\end{equation}

The Gaussian process prior is assumed to be separable:
\begin{equation}
f\sim \mathcal{GP}(0,k),
\qquad
k\big((s,t),(s',t')\big)
=
k_s(s,s')k_t(t,t').
\label{eq:separable-kernel}
\end{equation}
If latent values are vectorized over a full time--space grid, then
\begin{equation}
K_{ff}=K_t\otimes K_s.
\label{eq:full-kron-prior}
\end{equation}

The key modeling decision in this document is to train the linear mean and the GP
component jointly. Instead of first estimating $\beta$ and then fitting a GP to residuals,
we directly construct an online posterior over the joint unknowns $(\beta,u)$, where
$u$ denotes the GP inducing variables.

\section{Kronecker HiPPO-SVGP Inducing Representation}

\subsection{Temporal HiPPO-RFF Interdomain Features}

For the temporal axis, we use HiPPO-LegS interdomain inducing variables together with
a random Fourier feature approximation of a stationary temporal kernel.

For a time horizon $T>0$, define the scaled Legendre basis
\begin{equation}
g^{(T)}_\ell(x)
=
\sqrt{\frac{2\ell+1}{T}}\,
P_\ell\!\left(\frac{2x}{T}-1\right),
\qquad
x\in[0,T],
\quad
\ell=0,\dots,M_t-1,
\label{eq:scaled-legs-basis}
\end{equation}
where $P_\ell$ is the Legendre polynomial of degree $\ell$.
The temporal interdomain inducing variables are
\begin{equation}
u_\ell^{(t)}(T)
=
\int_0^T g_\ell^{(T)}(x) f_t(x)\,dx.
\label{eq:temporal-hippo-inducing}
\end{equation}

Assume the temporal kernel $k_t$ is stationary. By Bochner's theorem, we approximate it
using random Fourier features:
\begin{equation}
k_t(\tau,\tau')
\approx
\sum_{r=1}^{R} w_r
\left[
\cos(\omega_r\tau)\cos(\omega_r\tau')
+
\sin(\omega_r\tau)\sin(\omega_r\tau')
\right].
\label{eq:rff-kernel}
\end{equation}
Define the oscillatory integrals
\begin{align}
I^{\cos}_\ell(T;\omega)
&=
\int_0^T \cos(\omega x)g_\ell^{(T)}(x)\,dx,\\
I^{\sin}_\ell(T;\omega)
&=
\int_0^T \sin(\omega x)g_\ell^{(T)}(x)\,dx.
\end{align}
Using the spherical-Bessel identity
\begin{equation}
\int_{-1}^{1} e^{i\kappa u}P_\ell(u)\,du
=
2 i^\ell j_\ell(\kappa),
\qquad
\kappa=\frac{\omega T}{2},
\end{equation}
these integrals admit closed-form expressions involving spherical Bessel functions
$j_\ell(\cdot)$. This gives analytic temporal inducing covariances and point-to-inducing
cross-covariances.

Let $B_t(T)\in\mathbb{R}^{M_t\times 2R}$ be the temporal HiPPO-RFF feature matrix,
with entries
\begin{equation}
[B_t(T)]_{\ell,r}
=
\sqrt{w_r}\,I^{\cos}_\ell(T;\omega_r),
\qquad
[B_t(T)]_{\ell,R+r}
=
\sqrt{w_r}\,I^{\sin}_\ell(T;\omega_r),
\end{equation}
and define the temporal RFF vector
\begin{equation}
\varphi_t(\tau)
=
\left(
\sqrt{w_1}\cos(\omega_1\tau),\dots,\sqrt{w_R}\cos(\omega_R\tau),
\sqrt{w_1}\sin(\omega_1\tau),\dots,\sqrt{w_R}\sin(\omega_R\tau)
\right)^\top.
\end{equation}
Then
\begin{equation}
K_{uu}^{(t)}(T)
\approx
B_t(T)B_t(T)^\top,
\label{eq:temporal-inducing-cov}
\end{equation}
and
\begin{equation}
k_{fu}^{(t)}(\tau;T)
\approx
\varphi_t(\tau)^\top B_t(T)^\top.
\label{eq:temporal-cross-cov}
\end{equation}

For local temporal mini-batches, the same construction is applied using a local horizon
$T_n$ for block $D_n$, yielding $K_{uu,n}^{(t)}$ and $K_{fu,n}^{(t)}$.

\subsection{Spatial Inducing Variables}

On the spatial axis, we use inducing locations
\begin{equation}
Z_s=\{z_1^{(s)},\dots,z_{M_s}^{(s)}\}.
\end{equation}
The spatial inducing covariance and cross-covariance are
\begin{align}
K_{ZZ}^{(s)}
&\in\mathbb{R}^{M_s\times M_s},
&
[K_{ZZ}^{(s)}]_{jj'}
&=
k_s(z_j^{(s)},z_{j'}^{(s)}),
\\
K_{XZ}^{(s)}
&\in\mathbb{R}^{N_s\times M_s},
&
[K_{XZ}^{(s)}]_{ij}
&=
k_s(s_i,z_j^{(s)}).
\end{align}

The spatial inducing set is fixed across online time steps. This is important: online
adaptation occurs through the temporal HiPPO component, while the spatial inducing
geometry remains shared across all blocks.

\subsection{Mixed Spatio-Temporal Inducing Variables}

For each spatial inducing location $z_j^{(s)}$, define temporal interdomain variables
\begin{equation}
u_{\ell,j}(T)
=
\int_0^T g_\ell^{(T)}(x)f(x,z_j^{(s)})\,dx.
\end{equation}
Stacking the temporal and spatial inducing variables gives
\begin{equation}
u(T)
=
\vecop\!\left([u_{\ell,j}(T)]_{\ell,j}\right)
\in\mathbb{R}^{M_tM_s}.
\end{equation}

By separability of the spatio-temporal kernel,
\begin{equation}
K_{uu}^{(st)}(T)
=
K_{uu}^{(t)}(T)
\otimes
K_{ZZ}^{(s)}.
\label{eq:mixed-inducing-cov}
\end{equation}
For a block $D_n$, the point-to-inducing covariance has the corresponding Kronecker
form
\begin{equation}
K_{fu,n}^{(st)}
=
K_{fu,n}^{(t)}
\otimes
K_{XZ}^{(s)}.
\label{eq:mixed-cross-cov}
\end{equation}

The sparse GP projection matrix for block $n$ is therefore
\begin{equation}
A_n
=
K_{fu,n}^{(st)}
\left(K_{uu,n}^{(st)}\right)^{-1}.
\label{eq:projection-general}
\end{equation}
Using Kronecker identities,
\begin{equation}
A_n
=
\left[
K_{fu,n}^{(t)}
\left(K_{uu,n}^{(t)}\right)^{-1}
\right]
\otimes
\left[
K_{XZ}^{(s)}
\left(K_{ZZ}^{(s)}\right)^{-1}
\right].
\label{eq:projection-kron}
\end{equation}
Define
\begin{equation}
T_n
:=
K_{fu,n}^{(t)}
\left(K_{uu,n}^{(t)}\right)^{-1},
\qquad
C
:=
K_{XZ}^{(s)}
\left(K_{ZZ}^{(s)}\right)^{-1}.
\label{eq:Tn-C-def}
\end{equation}
Then
\begin{equation}
A_n=T_n\otimes C.
\label{eq:A-kron}
\end{equation}

\section{Why Residual-Based Online Learning is Inconsistent}

A two-stage procedure first estimates the linear mean and then fits a GP to residuals:
\begin{equation}
r_n(\hat{\beta}_n)
=
y_n-\Phi_n\hat{\beta}_n.
\end{equation}
The difficulty is that residuals are not fixed targets. If the estimate of the linear
coefficient changes from $\hat{\beta}_{n-1}$ to $\hat{\beta}_{n}$, then for any historical
block $j<n$,
\begin{align}
r_j(\hat{\beta}_{n})
&=
y_j-\Phi_j\hat{\beta}_{n}
\\
&=
y_j-\Phi_j\hat{\beta}_{n-1}
-
\Phi_j(\hat{\beta}_{n}-\hat{\beta}_{n-1})
\\
&=
r_j(\hat{\beta}_{n-1})
-
\Phi_j(\hat{\beta}_{n}-\hat{\beta}_{n-1}).
\label{eq:residual-mismatch}
\end{align}
Thus, the historical residual memory becomes stale whenever $\beta$ is updated.

\begin{proposition}[Joint training removes stale residual targets]
In the joint model
\begin{equation}
y_j\mid \beta,u_j
\sim
\mathcal{N}(\Phi_j\beta+A_j u_j,\sigma^2I),
\end{equation}
historical information is stored as a likelihood contribution to the same joint posterior
over $(\beta,u)$, rather than as residual observations computed from a changing point
estimate of $\beta$. Therefore, historical information does not become inconsistent when
the posterior mean of $\beta$ changes.
\end{proposition}

\begin{proof}
The two-stage target for block $j$ is
$r_j(\hat{\beta})=y_j-\Phi_j\hat{\beta}$, which explicitly depends on the current point
estimate $\hat{\beta}$. Equation~\eqref{eq:residual-mismatch} shows that the target
changes when $\hat{\beta}$ changes.

By contrast, joint training uses the likelihood
\begin{equation}
p(y_j\mid \beta,u_j)
=
\mathcal{N}(y_j\mid \Phi_j\beta+A_j u_j,\sigma^2I).
\end{equation}
This likelihood is a fixed probabilistic statement about $y_j$ as a function of the
unknown variables $(\beta,u_j)$. Updating the posterior mean of $\beta$ changes our
belief about $\beta$, but it does not redefine the historical observation $y_j$ or the
likelihood contribution of block $j$. Hence no stale residual target is created.
\end{proof}

The essential distinction is
\begin{equation}
\text{two-stage: learn } y-\Phi\hat{\beta},
\qquad
\text{joint training: learn } p(y\mid \beta,u).
\end{equation}
This is the reason the rest of this document focuses on a joint online posterior over
the linear mean and GP inducing variables.

\section{Online Joint Training ELBO}

For online block $n$, the sparse GP approximation gives
\begin{equation}
f_n\mid u_n \approx A_nu_n.
\end{equation}
Thus the likelihood is
\begin{equation}
p(y_n\mid \beta,u_n)
=
\mathcal{N}
\left(
y_n\mid \Phi_n\beta+A_nu_n,\sigma^2I
\right).
\label{eq:joint-likelihood}
\end{equation}

Let the streaming priors entering block $n$ be
\begin{align}
p_n(\beta)
&=
\mathcal{N}
\left(
m_{\beta,n|n-1},
S_{\beta,n|n-1}
\right),
\\
p_n(u_n)
&=
\mathcal{N}
\left(
m_{u,n|n-1},
S_{u,n|n-1}
\right).
\end{align}
A general joint Gaussian variational posterior is
\begin{equation}
q_n(\beta,u_n)
=
\mathcal{N}
\left(
\begin{bmatrix}\beta\\ u_n\end{bmatrix}
\middle|
m_n,S_n
\right).
\label{eq:joint-posterior}
\end{equation}
The online joint training objective is
\begin{equation}
\mathcal{L}_n
=
\E_{q_n(\beta,u_n)}
\left[
\log p(y_n\mid \beta,u_n)
\right]
-
\KL\!\left[
q_n(\beta,u_n)
\,\Vert\,
p_n(\beta,u_n)
\right].
\label{eq:online-joint-elbo}
\end{equation}

A practical factorized version is
\begin{equation}
q_n(\beta,u_n)
=
q_n(\beta)q_n(u_n),
\end{equation}
with
\begin{equation}
q_n(\beta)=\mathcal{N}(m_\beta,S_\beta),
\qquad
q_n(u_n)=\mathcal{N}(m_u,S_u).
\end{equation}
If the streaming prior also factorizes,
\begin{equation}
p_n(\beta,u_n)=p_n(\beta)p_n(u_n),
\end{equation}
then
\begin{align}
\mathcal{L}_n
&=
\E_{q_n(\beta)q_n(u_n)}
\left[
\log p(y_n\mid \beta,u_n)
\right]
\\
&\quad
-
\KL[q_n(\beta)\Vert p_n(\beta)]
-
\KL[q_n(u_n)\Vert p_n(u_n)].
\label{eq:mean-field-elbo}
\end{align}
For the Gaussian likelihood,
\begin{align}
\E\log p(y_n\mid \beta,u_n)
&=
-\frac{B_n}{2}\log(2\pi\sigma^2)
\\
&\quad
-\frac{1}{2\sigma^2}
\left[
\|y_n-\Phi_nm_\beta-A_nm_u\|^2
+
\tr(\Phi_nS_\beta\Phi_n^\top)
+
\tr(A_nS_uA_n^\top)
\right].
\label{eq:expected-log-likelihood}
\end{align}

This ELBO is the formal bridge between the joint linear-mean model and sparse
variational GP training. When the inducing coordinates are fixed, the previous posterior
can be used directly as the next streaming prior. When the HiPPO inducing coordinates
change, however, the expression \(p_n(u_n)\) must be constructed carefully. The main
method in this document uses an SSGP-style old-likelihood-ratio transfer, derived in
Section~\ref{sec:changing_basis_transfer}. The simpler prior-based ELBO remains useful
as a baseline form and as the projected-prior approximation discussed later.

\section{Fixed-Basis Exact Joint Gaussian Update}

If the inducing representation is fixed across online blocks, i.e. the same inducing
variable $u$ is used throughout, then the joint model is conjugate under a Gaussian
likelihood. In this special case, the exact joint posterior can be maintained in
information form.

Define the joint state
\begin{equation}
z=
\begin{bmatrix}
\beta\\
u
\end{bmatrix},
\qquad
H_n=
\begin{bmatrix}
\Phi_n & A_n
\end{bmatrix}.
\label{eq:joint-state-fixed-basis}
\end{equation}
Then
\begin{equation}
y_n\mid z
\sim
\mathcal{N}(H_nz,\sigma^2I).
\end{equation}
Let
\begin{equation}
q_{n-1}(z)=\mathcal{N}(m_{n-1},S_{n-1}),
\end{equation}
and define the natural parameters
\begin{equation}
\Lambda_{n-1}=S_{n-1}^{-1},
\qquad
h_{n-1}=S_{n-1}^{-1}m_{n-1}.
\end{equation}
The exact Gaussian-conjugate update is
\begin{equation}
\Lambda_n
=
\Lambda_{n-1}
+
\frac{1}{\sigma^2}H_n^\top H_n,
\label{eq:joint-precision-update}
\end{equation}
\begin{equation}
h_n
=
h_{n-1}
+
\frac{1}{\sigma^2}H_n^\top y_n.
\label{eq:joint-info-update}
\end{equation}
The posterior moments are recovered by
\begin{equation}
S_n=\Lambda_n^{-1},
\qquad
m_n=S_nh_n.
\label{eq:joint-moment-recovery}
\end{equation}

Expanding the precision contribution gives
\begin{equation}
H_n^\top H_n
=
\begin{bmatrix}
\Phi_n^\top\Phi_n
&
\Phi_n^\top A_n
\\
A_n^\top\Phi_n
&
A_n^\top A_n
\end{bmatrix}.
\label{eq:joint-precision-blocks}
\end{equation}
Therefore,
\begin{equation}
\Lambda_n
=
\Lambda_{n-1}
+
\frac{1}{\sigma^2}
\begin{bmatrix}
\Phi_n^\top\Phi_n
&
\Phi_n^\top A_n
\\
A_n^\top\Phi_n
&
A_n^\top A_n
\end{bmatrix},
\end{equation}
and
\begin{equation}
h_n
=
h_{n-1}
+
\frac{1}{\sigma^2}
\begin{bmatrix}
\Phi_n^\top y_n\\
A_n^\top y_n
\end{bmatrix}.
\end{equation}

This update is stronger than a mean-field approximation because it keeps the posterior
cross-covariance
\begin{equation}
\cov(\beta,u)\neq 0.
\end{equation}
The fixed-basis exact joint Gaussian update should therefore be used as a baseline and
sanity check for more scalable structured or mean-field approximations.

\begin{remark}
Equations~\eqref{eq:joint-precision-update}--\eqref{eq:joint-moment-recovery} are valid
only when all historical likelihood contributions are expressed in the same inducing
coordinates. If the HiPPO temporal inducing basis changes over time, the previous
posterior must first be transferred to the new inducing coordinates before the next
likelihood update is applied.
\end{remark}

\section{Changing-Basis Posterior-to-Prior Transfer}
\label{sec:changing_basis_transfer}

In the HiPPO-SVGP setting, the temporal inducing basis depends on the current horizon
or on the current local temporal block. Consequently, the inducing variables used at two
consecutive online steps are generally not the same random variables. Let \(u_o\) denote
the old inducing variable at step \(n-1\), and let \(u_n\) denote the new inducing
variable at step \(n\). In general,
\begin{equation}
    u_o \neq u_n .
\end{equation}
For example, the old and new temporal HiPPO interdomain variables may be defined as
\begin{equation}
    u_{\ell}^{old}
    =
    \int_0^{T_{n-1}}
    g_\ell^{(T_{n-1})}(x)f(x)\,dx,
    \qquad
    u_{\ell}^{new}
    =
    \int_0^{T_n}
    g_\ell^{(T_n)}(x)f(x)\,dx .
\end{equation}
Because the basis function \(g_\ell^{(T)}\) changes with the horizon \(T\), the old
posterior \(q_{n-1}(u_o)\) cannot be directly reused as a posterior or prior over
\(u_n\). Therefore, the naive assignment
\begin{equation}
    p_n(u_n)=q_{n-1}(u_o)
\end{equation}
is not well-defined.

The central problem is thus to transfer the information contained in \(q_{n-1}(u_o)\)
into the new inducing coordinates \(u_n\). We distinguish two related transfer
mechanisms. The first one is a streaming variational free-energy transfer that explicitly
propagates the old likelihood information, following the principle of streaming sparse
Gaussian processes. This is the main theoretical formulation. The second one is a simpler
Gaussian projected-prior transfer, which is useful as an efficient approximation and as an
ablation baseline.

\subsection{Streaming VFE Transfer with Changing HiPPO Coordinates}
\label{subsec:streaming_vfe_transfer}

The streaming sparse GP framework represents old observations without replaying them.
The old data contribution is encoded by the ratio between the previous posterior and the
old prior. In the standard notation, let \(a\) denote the old inducing variables and let
\(b\) denote the new inducing variables. The old data contribution is represented through
\begin{equation}
    \frac{q_{\mathrm{old}}(a)}{p(a)} .
\end{equation}
The corresponding variational free-energy update can be written as
\begin{equation}
    q_{\mathrm{vfe}}(b)
    \propto
    p(b)
    \exp
    \left\{
    \E_{p(a\mid b)}
    \left[
        \log
        \frac{q_{\mathrm{old}}(a)}{p(a)}
    \right]
    +
    \E_{p(f_{\mathrm{new}}\mid b)}
    \left[
        \log p(y_{\mathrm{new}}\mid f_{\mathrm{new}})
    \right]
    \right\}.
    \label{eq:ssgp_vfe_original}
\end{equation}
The first expectation transfers the old approximate likelihood factor into the new
inducing representation, while the second expectation assimilates the likelihood
contribution of the new data block.

In our model, the same principle is applied to HiPPO interdomain inducing variables.
The old inducing variable is \(u_o\), the new inducing variable is \(u_n\), and the new
observation model is the joint linear-mean GP likelihood
\begin{equation}
    p(y_n\mid \beta,u_n)
    =
    \mathcal{N}
    \left(
        y_n
        \mid
        \Phi_n\beta + A_nu_n,
        \sigma^2 I
    \right).
    \label{eq:joint_likelihood_changing_basis}
\end{equation}
The direct analogue of the streaming sparse GP VFE update is therefore
\begin{equation}
    q_n(u_n)
    \propto
    p(u_n)
    \exp
    \left\{
    \E_{p(u_o\mid u_n)}
    \left[
        \log
        \frac{
            q_{n-1}(u_o)
        }{
            p(u_o)
        }
    \right]
    +
    \E_{p(f_n\mid u_n)}
    \left[
        \log p(y_n\mid \beta,u_n)
    \right]
    \right\}.
    \label{eq:hippo_streaming_vfe_analogue}
\end{equation}
The first term transfers the approximate old likelihood information from the old HiPPO
basis into the new HiPPO basis. The second term incorporates the new data block under
the joint linear-mean GP likelihood.

Equivalently, define the old-information transfer term
\begin{equation}
    \mathcal{T}_{o\rightarrow n}(u_n)
    :=
    \E_{p(u_o\mid u_n)}
    \left[
        \log
        \frac{
            q_{n-1}(u_o)
        }{
            p(u_o)
        }
    \right].
    \label{eq:old_likelihood_transfer_term}
\end{equation}
Then
\begin{equation}
    q_n(u_n)
    \propto
    p(u_n)
    \exp
    \left\{
        \mathcal{T}_{o\rightarrow n}(u_n)
        +
        \E_{p(f_n\mid u_n)}
        \left[
            \log p(y_n\mid \beta,u_n)
        \right]
    \right\}.
\end{equation}
This formulation is stronger than a simple posterior projection because it propagates
the old likelihood factor \(q_{n-1}(u_o)/p(u_o)\), rather than only the marginal
predictive distribution of \(u_n\) implied by \(q_{n-1}(u_o)\).

\subsection{Closed-Form Gaussian Old-Likelihood Transfer}
\label{subsec:closed_form_old_likelihood_transfer}

We now derive the closed-form Gaussian version of
\eqref{eq:hippo_streaming_vfe_analogue}. Assume that the old posterior and old prior are
\begin{equation}
    q_{n-1}(u_o)
    =
    \mathcal{N}(u_o\mid m_o,S_o),
    \qquad
    p(u_o)
    =
    \mathcal{N}(u_o\mid 0,K_{oo}).
\end{equation}
Then the log old-likelihood ratio is quadratic:
\begin{align}
    \log
    \frac{
        q_{n-1}(u_o)
    }{
        p(u_o)
    }
    &=
    -\frac{1}{2}
    u_o^\top
    \left(
        S_o^{-1}-K_{oo}^{-1}
    \right)
    u_o
    +
    u_o^\top S_o^{-1}m_o
    +
    c_o ,
\end{align}
where \(c_o\) collects terms that do not depend on \(u_o\). Define the old likelihood
natural parameters
\begin{equation}
    R_o
    :=
    S_o^{-1}-K_{oo}^{-1},
    \qquad
    r_o
    :=
    S_o^{-1}m_o .
    \label{eq:old_likelihood_natural_params}
\end{equation}
Thus,
\begin{equation}
    \log
    \frac{
        q_{n-1}(u_o)
    }{
        p(u_o)
    }
    =
    -\frac{1}{2}u_o^\top R_o u_o
    +
    u_o^\top r_o
    +
    c_o .
    \label{eq:old_likelihood_ratio_quadratic}
\end{equation}

Under the GP prior, define the joint covariance between the new and old inducing
variables by
\begin{equation}
    \begin{bmatrix}
        u_n \\
        u_o
    \end{bmatrix}
    \sim
    \mathcal{N}
    \left(
    0,
    \begin{bmatrix}
        K_{nn} & K_{no} \\
        K_{on} & K_{oo}
    \end{bmatrix}
    \right).
    \label{eq:new_old_joint_prior}
\end{equation}
The conditional distribution of the old inducing variable given the new one is
\begin{equation}
    p(u_o\mid u_n)
    =
    \mathcal{N}
    \left(
        L_{on}u_n,
        \Sigma_{o\mid n}
    \right),
    \label{eq:old_given_new_conditional}
\end{equation}
where
\begin{equation}
    L_{on}
    :=
    K_{on}K_{nn}^{-1},
    \qquad
    \Sigma_{o\mid n}
    :=
    K_{oo}-K_{on}K_{nn}^{-1}K_{no}.
    \label{eq:old_given_new_params}
\end{equation}
Substituting \eqref{eq:old_likelihood_ratio_quadratic} into
\eqref{eq:old_likelihood_transfer_term} gives
\begin{align}
    \mathcal{T}_{o\rightarrow n}(u_n)
    &=
    \E_{p(u_o\mid u_n)}
    \left[
        -\frac{1}{2}u_o^\top R_o u_o
        +
        u_o^\top r_o
    \right]
    +
    \mathrm{const}  \\
    &=
    -\frac{1}{2}
    u_n^\top
    L_{on}^\top R_o L_{on}
    u_n
    +
    u_n^\top
    L_{on}^\top r_o
    +
    \mathrm{const}.
    \label{eq:transferred_old_likelihood_quadratic}
\end{align}
The trace term \(\tr(R_o\Sigma_{o\mid n})\) is absorbed into the constant because it
does not depend on \(u_n\).

Therefore, the old data contribute the following natural parameters in the new inducing
coordinates:
\begin{equation}
    \Lambda_{o\rightarrow n}
    =
    L_{on}^\top R_o L_{on},
    \qquad
    h_{o\rightarrow n}
    =
    L_{on}^\top r_o .
    \label{eq:old_to_new_natural_params}
\end{equation}
These are the old-likelihood-ratio contributions transferred from the old HiPPO basis to
the new HiPPO basis.

For the new data block, under the Gaussian observation model
\eqref{eq:joint_likelihood_changing_basis}, the GP part contributes
\begin{equation}
    \Lambda_{\mathrm{new},n}
    =
    \frac{1}{\sigma^2}A_n^\top A_n,
    \qquad
    h_{\mathrm{new},n}
    =
    \frac{1}{\sigma^2}
    A_n^\top
    \left(
        y_n-\Phi_n m_{\beta,n}
    \right)
\end{equation}
when using a mean-field update in which \(q_n(\beta)\) and \(q_n(u_n)\) are updated
alternately.

Combining the GP prior, the transferred old likelihood, and the new likelihood gives
\begin{equation}
    S_{u,n}^{-1}
    =
    K_{nn}^{-1}
    +
    \Lambda_{o\rightarrow n}
    +
    \frac{1}{\sigma^2}A_n^\top A_n,
    \label{eq:ssgp_style_u_precision}
\end{equation}
and
\begin{equation}
    S_{u,n}^{-1}m_{u,n}
    =
    h_{o\rightarrow n}
    +
    \frac{1}{\sigma^2}
    A_n^\top
    \left(
        y_n-\Phi_n m_{\beta,n}
    \right).
    \label{eq:ssgp_style_u_info_vector}
\end{equation}
Equations~\eqref{eq:ssgp_style_u_precision}--\eqref{eq:ssgp_style_u_info_vector}
are the SSGP-style changing-basis update for the HiPPO inducing variables. They
explicitly propagate old likelihood information through the ratio
\(q_{n-1}(u_o)/p(u_o)\), and therefore more closely follow the streaming sparse GP
variational principle than a simple posterior moment projection.

\subsection{Joint Linear-Mean Update}
\label{subsec:joint_linear_mean_update_transfer}

The linear coefficient \(\beta\) is updated jointly with the GP inducing state. If \(\beta\)
is assumed static, its streaming prior is simply
\begin{equation}
    p_n(\beta)=q_{n-1}(\beta).
\end{equation}
If slowly time-varying linear effects are allowed, we use a drift prior
\begin{equation}
    p_n(\beta)
    =
    \mathcal{N}
    \left(
        m_{\beta,n-1},
        S_{\beta,n-1}+Q_\beta
    \right),
    \qquad
    Q_\beta\succeq 0 .
    \label{eq:beta-drift-prior}
\end{equation}

For a mean-field posterior
\begin{equation}
    q_n(\beta,u_n)=q_n(\beta)q_n(u_n),
\end{equation}
the coordinate update for \(\beta\) is
\begin{equation}
    S_{\beta,n}^{-1}
    =
    S_{\beta,n\mid n-1}^{-1}
    +
    \frac{1}{\sigma^2}\Phi_n^\top\Phi_n,
    \label{eq:beta_coordinate_precision_transfer}
\end{equation}
and
\begin{equation}
    S_{\beta,n}^{-1}m_{\beta,n}
    =
    S_{\beta,n\mid n-1}^{-1}m_{\beta,n\mid n-1}
    +
    \frac{1}{\sigma^2}
    \Phi_n^\top
    \left(
        y_n-A_nm_{u,n}
    \right).
    \label{eq:beta_coordinate_info_transfer}
\end{equation}
The update for \(u_n\) is given by
\eqref{eq:ssgp_style_u_precision}--\eqref{eq:ssgp_style_u_info_vector}. In practice,
one alternates
\eqref{eq:beta_coordinate_precision_transfer}--\eqref{eq:beta_coordinate_info_transfer}
and \eqref{eq:ssgp_style_u_precision}--\eqref{eq:ssgp_style_u_info_vector} for a
small number of inner iterations per block.

This joint update is different from a residual-based two-stage method. The likelihood is
always applied to the original observations \(y_n\),
\begin{equation}
    p(y_n\mid \beta,u_n)
    =
    \mathcal{N}
    \left(
        y_n
        \mid
        \Phi_n\beta + A_nu_n,
        \sigma^2I
    \right),
\end{equation}
rather than to residuals computed using a changing point estimate of \(\beta\). Hence
historical information is carried through the joint posterior and through the
old-likelihood-ratio transfer, not through stale residual targets.

\subsection{Projected-Prior Approximation}
\label{subsec:projected_prior_approximation}

The SSGP-style transfer above is the main theoretical update. However, it requires
constructing and propagating the old-likelihood ratio
\begin{equation}
    \log
    \frac{q_{n-1}(u_o)}{p(u_o)}.
\end{equation}
For a simpler and cheaper implementation, one may instead use a Gaussian projected-prior
approximation. This approximation directly maps the old posterior over \(u_o\) into a
predictive Gaussian prior over \(u_n\):
\begin{equation}
    p_n^{\mathrm{proj}}(u_n)
    =
    \int
    p(u_n\mid u_o)
    q_{n-1}(u_o)
    \,du_o .
    \label{eq:posterior-to-prior-integral}
\end{equation}
Assume again that
\begin{equation}
    q_{n-1}(u_o)=\mathcal{N}(m_o,S_o),
\end{equation}
and use the GP conditional
\begin{equation}
    p(u_n\mid u_o)
    =
    \mathcal{N}
    \left(
        K_{no}K_{oo}^{-1}u_o,
        K_{nn}-K_{no}K_{oo}^{-1}K_{on}
    \right).
\end{equation}
Then
\begin{equation}
    p_n^{\mathrm{proj}}(u_n)
    =
    \mathcal{N}
    \left(
        m_{u,n\mid n-1}^{\mathrm{proj}},
        S_{u,n\mid n-1}^{\mathrm{proj}}
    \right),
\end{equation}
where
\begin{equation}
    m_{u,n\mid n-1}^{\mathrm{proj}}
    =
    K_{no}K_{oo}^{-1}m_o,
    \label{eq:transfer-mean}
\end{equation}
and
\begin{equation}
    S_{u,n\mid n-1}^{\mathrm{proj}}
    =
    K_{nn}
    +
    K_{no}K_{oo}^{-1}
    (S_o-K_{oo})
    K_{oo}^{-1}K_{on}.
    \label{eq:transfer-cov}
\end{equation}
This projected-prior update is a filtering-style moment projection. It preserves the
posterior mean and covariance information of the old HiPPO state and maps them into the
current basis. If the old posterior equals the old prior, \(S_o=K_{oo}\), then
\begin{equation}
    S_{u,n\mid n-1}^{\mathrm{proj}}=K_{nn},
\end{equation}
so no information beyond the GP prior is carried forward. Conversely, if \(S_o\neq K_{oo}\),
then the term \(S_o-K_{oo}\) transfers the posterior uncertainty reduction induced by
past data.

The projected-prior approximation is easier to implement and is useful as a computationally
efficient baseline. However, it does not explicitly propagate the old likelihood ratio
\(q_{n-1}(u_o)/p(u_o)\). Therefore, the SSGP-style transfer in
\eqref{eq:ssgp_style_u_precision}--\eqref{eq:ssgp_style_u_info_vector} is the preferred
formulation when the goal is to obtain a more principled streaming variational method.

\subsection{Summary of the Changing-Basis Transfer}
\label{subsec:changing_basis_summary}

The changing-basis problem arises because the HiPPO inducing coordinates evolve with the
temporal horizon or local block. The old inducing state \(u_o\) and the new inducing
state \(u_n\) are distinct random variables. This section introduced two transfer
mechanisms:
\begin{enumerate}[label=(\arabic*)]
    \item \textbf{SSGP-style old-likelihood-ratio transfer.} This is the main
    theoretical update. It propagates the old approximate likelihood factor
    \(q_{n-1}(u_o)/p(u_o)\) through \(p(u_o\mid u_n)\), giving the transferred natural
    parameters
    \begin{equation}
        \Lambda_{o\rightarrow n}
        =
        L_{on}^\top R_oL_{on},
        \qquad
        h_{o\rightarrow n}
        =
        L_{on}^\top r_o .
    \end{equation}
    The resulting posterior update is
    \begin{equation}
        S_{u,n}^{-1}
        =
        K_{nn}^{-1}
        +
        \Lambda_{o\rightarrow n}
        +
        \frac{1}{\sigma^2}A_n^\top A_n .
    \end{equation}

    \item \textbf{Gaussian projected-prior transfer.} This is a simpler filtering-style
    approximation:
    \begin{equation}
        p_n^{\mathrm{proj}}(u_n)
        =
        \int
        p(u_n\mid u_o)q_{n-1}(u_o)\,du_o .
    \end{equation}
    It is computationally convenient but less directly tied to the streaming sparse GP
    variational free-energy derivation.
\end{enumerate}

In the remainder of the method, the SSGP-style transfer is treated as the main principled
changing-basis update, while the projected-prior transfer serves as an efficient
approximation and an ablation baseline. This separation allows the model to retain a
strong connection to streaming variational sparse GP inference while still supporting
practical Kronecker-structured HiPPO-SVGP implementation.

\section{Mean-Field Coordinate Updates}
\label{sec:mean_field_coordinate_updates}

The full joint posterior over \((\beta,u_n)\) is accurate but may be expensive when
\(M=M_tM_s\) is large. A scalable engineering version is the factorized posterior
\begin{equation}
    q_n(\beta,u_n)=q_n(\beta)q_n(u_n).
\end{equation}
This section gives the coordinate updates consistent with the SSGP-style transfer in
Section~\ref{sec:changing_basis_transfer}. The projected-prior approximation is also
included as an implementation baseline.

\subsection{Coordinate Update for the Linear Coefficient}
\label{subsec:coordinate_beta}

Let the streaming prior for \(\beta\) entering block \(n\) be
\begin{equation}
    p_n(\beta)
    =
    \mathcal{N}
    \left(
        m_{\beta,n\mid n-1},
        S_{\beta,n\mid n-1}
    \right).
\end{equation}
Holding \(q_n(u_n)\) fixed, the coordinate-optimal Gaussian factor for \(\beta\) is
\begin{equation}
    S_{\beta,n}^{-1}
    =
    S_{\beta,n\mid n-1}^{-1}
    +
    \frac{1}{\sigma^2}\Phi_n^\top\Phi_n,
    \label{eq:mf_beta_precision}
\end{equation}
and
\begin{equation}
    S_{\beta,n}^{-1}m_{\beta,n}
    =
    S_{\beta,n\mid n-1}^{-1}m_{\beta,n\mid n-1}
    +
    \frac{1}{\sigma^2}
    \Phi_n^\top
    \left(
        y_n-A_nm_{u,n}
    \right).
    \label{eq:mf_beta_info}
\end{equation}
This resembles a residual update, but it is only an inner coordinate step of one joint
ELBO. No historical residual target is stored.

\subsection{Coordinate Update for the GP Inducing State under SSGP-Style Transfer}
\label{subsec:coordinate_u_ssgp_transfer}

Holding \(q_n(\beta)\) fixed, the coordinate-optimal Gaussian factor for \(u_n\) uses
three sources of information: the current GP prior \(p(u_n)=\mathcal{N}(0,K_{nn})\),
the transferred old-likelihood factor, and the new block likelihood. Therefore,
\begin{equation}
    S_{u,n}^{-1}
    =
    K_{nn}^{-1}
    +
    \Lambda_{o\rightarrow n}
    +
    \frac{1}{\sigma^2}A_n^\top A_n,
    \label{eq:mf_u_precision_ssgp_transfer}
\end{equation}
and
\begin{equation}
    S_{u,n}^{-1}m_{u,n}
    =
    h_{o\rightarrow n}
    +
    \frac{1}{\sigma^2}
    A_n^\top
    \left(
        y_n-\Phi_n m_{\beta,n}
    \right).
    \label{eq:mf_u_info_ssgp_transfer}
\end{equation}
Here
\begin{equation}
    \Lambda_{o\rightarrow n}
    =
    L_{on}^\top R_oL_{on},
    \qquad
    h_{o\rightarrow n}
    =
    L_{on}^\top r_o,
\end{equation}
with
\begin{equation}
    R_o=S_o^{-1}-K_{oo}^{-1},
    \qquad
    r_o=S_o^{-1}m_o.
\end{equation}
Equations~\eqref{eq:mf_u_precision_ssgp_transfer}--\eqref{eq:mf_u_info_ssgp_transfer}
are the mean-field coordinate update corresponding to the SSGP-style
old-likelihood-ratio transfer.

In practice, one performs a small number of inner sweeps per online block:
\begin{enumerate}[label=(\roman*)]
    \item initialize \(q_n(\beta)\) from \(p_n(\beta)\);
    \item construct the old-likelihood transfer
    \((\Lambda_{o\rightarrow n},h_{o\rightarrow n})\);
    \item update \(q_n(u_n)\) using
    \eqref{eq:mf_u_precision_ssgp_transfer}--\eqref{eq:mf_u_info_ssgp_transfer};
    \item update \(q_n(\beta)\) using
    \eqref{eq:mf_beta_precision}--\eqref{eq:mf_beta_info};
    \item repeat the last two steps if needed.
\end{enumerate}

\subsection{Projected-Prior Coordinate Update}
\label{subsec:coordinate_projected_prior}

For comparison, the projected-prior approximation replaces the transferred
old-likelihood factor by a Gaussian streaming prior
\begin{equation}
    p_n^{\mathrm{proj}}(u_n)
    =
    \mathcal{N}
    \left(
        m_{u,n\mid n-1}^{\mathrm{proj}},
        S_{u,n\mid n-1}^{\mathrm{proj}}
    \right).
\end{equation}
In that approximation,
\begin{equation}
    S_{u,n}^{-1}
    =
    \left(
        S_{u,n\mid n-1}^{\mathrm{proj}}
    \right)^{-1}
    +
    \frac{1}{\sigma^2}A_n^\top A_n,
\end{equation}
and
\begin{equation}
    S_{u,n}^{-1}m_{u,n}
    =
    \left(
        S_{u,n\mid n-1}^{\mathrm{proj}}
    \right)^{-1}
    m_{u,n\mid n-1}^{\mathrm{proj}}
    +
    \frac{1}{\sigma^2}
    A_n^\top
    \left(
        y_n-\Phi_nm_{\beta,n}
    \right).
\end{equation}
This projected-prior version is used as an efficient approximation or ablation baseline,
not as the main theoretical update.

\section{Kronecker-Aware Implementation}
\label{sec:kronecker_aware_implementation}

The Kronecker form
\begin{equation}
    A_n=T_n\otimes C
\end{equation}
should be used explicitly. Therefore,
\begin{equation}
    A_n^\top A_n
    =
    (T_n^\top T_n)\otimes(C^\top C).
\end{equation}
Let
\begin{equation}
    G:=C^\top C.
\end{equation}
Since the spatial inducing set is fixed, \(G\) can be precomputed.

Under the SSGP-style changing-basis transfer, the GP inducing precision contains three
terms:
\begin{equation}
    \Lambda_{u,n}
    =
    K_{nn}^{-1}
    +
    \Lambda_{o\rightarrow n}
    +
    \frac{1}{\sigma^2}
    A_n^\top A_n.
    \label{eq:kronecker_precision_with_old_transfer}
\end{equation}
The main challenge is the transferred old-likelihood precision
\(\Lambda_{o\rightarrow n}=L_{on}^\top R_oL_{on}\). In a generic model this term need
not be Kronecker-separable. In the present construction, however, it is
Kronecker-preserving because the inducing basis changes only along the temporal
dimension, while the spatial inducing geometry remains fixed.

\subsection{Kronecker-Preserving Old-Likelihood Transfer}
\label{subsec:kronecker_preserving_transfer}

Let
\begin{equation}
    K_{nn}
    =
    K_{nn}^{(t)}\otimes K_s,
    \qquad
    K_{oo}
    =
    K_{oo}^{(t)}\otimes K_s,
    \qquad
    K_{on}
    =
    K_{on}^{(t)}\otimes K_s,
    \label{eq:old_new_kronecker_covariances}
\end{equation}
where \(K_s=K_{ZZ}^{(s)}\). Since the spatial inducing locations are fixed, the spatial
factor is the same in the old, new, and cross covariances. Then
\begin{align}
    L_{on}
    &=
    K_{on}K_{nn}^{-1} \\
    &=
    \left(
        K_{on}^{(t)}\otimes K_s
    \right)
    \left[
        \left(K_{nn}^{(t)}\right)^{-1}
        \otimes
        K_s^{-1}
    \right] \\
    &=
    \left[
        K_{on}^{(t)}
        \left(K_{nn}^{(t)}\right)^{-1}
    \right]
    \otimes I_s .
\end{align}
Define
\begin{equation}
    L_{on}^{(t)}
    :=
    K_{on}^{(t)}
    \left(K_{nn}^{(t)}\right)^{-1}.
\end{equation}
Then
\begin{equation}
    L_{on}=L_{on}^{(t)}\otimes I_s.
    \label{eq:temporal_only_transfer_operator}
\end{equation}

Moreover, because each block likelihood projection has the form
\begin{equation}
    A_j=T_j\otimes C,
\end{equation}
each Gaussian likelihood contribution to the inducing precision is
\begin{equation}
    \frac{1}{\sigma^2}A_j^\top A_j
    =
    \frac{1}{\sigma^2}
    (T_j^\top T_j)
    \otimes
    G.
\end{equation}
We therefore maintain the old likelihood precision directly in Kronecker natural form,
\begin{equation}
    R_o=B_o\otimes G,
    \label{eq:old_likelihood_precision_kronecker}
\end{equation}
rather than recovering it by dense subtraction \(R_o=S_o^{-1}-K_{oo}^{-1}\). This is
important: forming a dense posterior covariance and then subtracting the prior precision
may destroy the structure numerically, while the natural likelihood statistic remains
structured by construction.

Substituting \eqref{eq:temporal_only_transfer_operator} and
\eqref{eq:old_likelihood_precision_kronecker} into the old-likelihood transfer gives
\begin{align}
    \Lambda_{o\rightarrow n}
    &=
    L_{on}^\top R_oL_{on} \\
    &=
    \left[
        \left(L_{on}^{(t)}\right)^\top
        \otimes I_s
    \right]
    (B_o\otimes G)
    \left[
        L_{on}^{(t)}\otimes I_s
    \right] \\
    &=
    \left[
        \left(L_{on}^{(t)}\right)^\top
        B_o
        L_{on}^{(t)}
    \right]
    \otimes G.
\end{align}
Hence the transferred old likelihood remains Kronecker-separable. Define
\begin{equation}
    B_{o\rightarrow n}
    :=
    \left(L_{on}^{(t)}\right)^\top
    B_o
    L_{on}^{(t)}.
    \label{eq:transferred_temporal_likelihood_statistic}
\end{equation}
Then
\begin{equation}
    \Lambda_{o\rightarrow n}
    =
    B_{o\rightarrow n}\otimes G.
    \label{eq:kronecker_old_likelihood_transfer}
\end{equation}

The new block likelihood contributes
\begin{equation}
    \frac{1}{\sigma^2}A_n^\top A_n
    =
    \frac{1}{\sigma^2}
    (T_n^\top T_n)\otimes G.
\end{equation}
Therefore the temporal likelihood statistic is updated by
\begin{equation}
    B_n
    =
    B_{o\rightarrow n}
    +
    \frac{1}{\sigma^2}T_n^\top T_n.
    \label{eq:temporal_stat_ssgp_transfer}
\end{equation}
The resulting inducing precision is
\begin{equation}
    \Lambda_{u,n}
    =
    \left(K_{nn}^{(t)}\right)^{-1}
    \otimes
    K_s^{-1}
    +
    B_n\otimes G.
    \label{eq:kronecker_precision_after_ssgp_transfer}
\end{equation}
Thus, the SSGP-style old-likelihood-ratio transfer does not destroy the Sylvester
structure, provided that the spatial inducing set is fixed and the changing basis acts
only along the temporal dimension.

\begin{remark}
If the spatial observation pattern or spatial projection changes across online blocks, then
the spatial factor \(G\) may become block-dependent. The precision then has the more
general form
\[
    K_t^{-1}\otimes K_s^{-1}
    +
    \sum_j B_j\otimes G_j,
\]
which no longer reduces to the single Sylvester equation below. It can still be handled by
Kronecker matrix-vector products and iterative solvers, but the closed-form Sylvester
diagonalization is exact only when the spatial factor is shared or approximated by a
shared structured factor.
\end{remark}

\subsection{Information Vector Transfer}
\label{subsec:information_vector_transfer}

The old information vector is transferred by
\begin{equation}
    h_{o\rightarrow n}
    =
    L_{on}^\top r_o
    =
    \left[
        \left(L_{on}^{(t)}\right)^\top
        \otimes I_s
    \right]r_o.
    \label{eq:kronecker_information_transfer}
\end{equation}
This multiplication should be implemented as a Kronecker matrix-vector product. If
\(r_o=\vecop(R_{\mathrm{info},o})\), the transformation corresponds to multiplying the
reshaped information matrix on the temporal side; the exact left/right convention depends
only on the chosen vectorization order.

The new block information vector is
\begin{equation}
    h_{\mathrm{new},n}
    =
    \frac{1}{\sigma^2}
    A_n^\top
    \left(
        y_n-\Phi_nm_{\beta,n}
    \right),
\end{equation}
and it should also be computed using \(A_n=T_n\otimes C\), without materializing the full
matrix \(A_n\).

\subsection{Sylvester Reformulation for Posterior Mean and Uncertainty}
\label{subsec:sylvester_reformulation}

With the Kronecker-preserving transfer in
\eqref{eq:kronecker_precision_after_ssgp_transfer}, linear systems involving
\(\Lambda_{u,n}\) can still be reduced to Sylvester equations. This applies both to the
posterior mean solve and to predictive uncertainty quadratic forms.

Let
\begin{equation}
    q=\vecop(Q),
    \qquad
    z=\Lambda_{u,n}^{-1}q=\vecop(Z),
\end{equation}
where \(Q,Z\in\mathbb{R}^{M_s\times M_t}\). Using
\begin{equation}
    \Lambda_{u,n}
    =
    \left(K_{nn}^{(t)}\right)^{-1}
    \otimes
    K_s^{-1}
    +
    B_n\otimes G,
\end{equation}
and the identity
\begin{equation}
    (A\otimes B)\vecop(Z)
    =
    \vecop(BZA^\top),
\end{equation}
the linear system
\begin{equation}
    \Lambda_{u,n}z=q
\end{equation}
is equivalent to
\begin{equation}
    K_s^{-1}Z
    \left(K_{nn}^{(t)}\right)^{-1}
    +
    GZB_n
    =
    Q,
    \label{eq:sylvester_after_ssgp_transfer}
\end{equation}
assuming \(K_{nn}^{(t)}\) and \(B_n\) are symmetric. After solving for \(Z\),
\begin{equation}
    q^\top\Lambda_{u,n}^{-1}q
    =
    \langle Q,Z\rangle_F.
\end{equation}
Thus, the same Sylvester mechanism can be used after SSGP-style transfer. The key is to
maintain the old likelihood precision as the temporal statistic \(B_o\) rather than as a
dense matrix.

\section{Prediction and Uncertainty}
\label{sec:prediction_uncertainty}

For a test point \(x_*=(t_*,s_*)\), let
\begin{equation}
    \phi_*=\phi(t_*,s_*),
\end{equation}
and define
\begin{equation}
    a_*^\top=K_{*u}K_{uu}^{-1}.
\end{equation}
In the main method, the posterior moments \(m_u\) and \(S_u\) used below are obtained
from the SSGP-style old-likelihood-ratio transfer update in
\eqref{eq:ssgp_style_u_precision}--\eqref{eq:ssgp_style_u_info_vector}. When the
projected-prior approximation is used as an ablation, the same predictive equations
apply, but \(m_u\) and \(S_u\) are produced by the projected-prior coordinate update.

Under the mean-field posterior \(q(\beta)q(u)\), the predictive mean is
\begin{equation}
    \E[y_*]
    =
    \phi_*^\top m_\beta
    +
    a_*^\top m_u .
\end{equation}
The predictive variance is
\begin{equation}
    \Var(y_*)
    =
    \sigma^2
    +
    \phi_*^\top S_\beta\phi_*
    +
    k(x_*,x_*)
    -
    K_{*u}K_{uu}^{-1}K_{u*}
    +
    a_*^\top S_ua_* .
    \label{eq:mean_field_predictive_variance}
\end{equation}
The five terms correspond to observation noise, linear-mean posterior uncertainty, prior
GP variance, inducing conditioning variance reduction, and posterior uncertainty in the
GP inducing variables. The last term can be evaluated by the Sylvester solve in
Section~\ref{subsec:sylvester_reformulation}.

Under a full joint posterior
\begin{equation}
    q(z)
    =
    \mathcal{N}(m_z,S_z),
    \qquad
    z=
    \begin{bmatrix}
        \beta\\
        u
    \end{bmatrix},
\end{equation}
define
\begin{equation}
    h_*=
    \begin{bmatrix}
        \phi_*\\
        a_*
    \end{bmatrix}.
\end{equation}
The predictive mean is
\begin{equation}
    \E[y_*]=h_*^\top m_z,
\end{equation}
and the predictive variance is
\begin{equation}
    \Var(y_*)
    =
    \sigma^2
    +
    \Var(f_*\mid u)
    +
    h_*^\top S_zh_* .
\end{equation}
If
\begin{equation}
    S_z=
    \begin{bmatrix}
        S_{\beta\beta} & S_{\beta u}\\
        S_{u\beta} & S_{uu}
    \end{bmatrix},
\end{equation}
then
\begin{equation}
    h_*^\top S_zh_*
    =
    \phi_*^\top S_{\beta\beta}\phi_*
    +
    a_*^\top S_{uu}a_*
    +
    2\phi_*^\top S_{\beta u}a_* .
\end{equation}
The cross term \(2\phi_*^\top S_{\beta u}a_*\) is absent in the mean-field
approximation. This is why the fixed-basis exact joint Gaussian update remains useful as
a calibration baseline.

\section{Experiments and Ablations}
\label{sec:experiments_ablations}

The recommended experimental sequence is designed to isolate the contribution of each
modeling choice: joint training, changing-basis transfer, old-likelihood-ratio transfer,
Kronecker structure, and posterior factorization.

\subsection{Methods to Compare}

At minimum, compare the following methods:
\begin{enumerate}[label=(\arabic*)]
    \item \textbf{Two-stage future-only baseline.}
    The linear mean is updated first and the GP is trained on current residuals only.

    \item \textbf{Two-stage residual GP with buffer or sketch correction.}
    This tests whether explicit residual repair can mitigate stale residual targets.

    \item \textbf{Fixed-basis exact joint Gaussian baseline.}
    This uses a fixed inducing representation and performs the exact joint Gaussian
    update over \((\beta,u)\).

    \item \textbf{Joint mean-field without changing-basis transfer.}
    This tests the failure mode caused by ignoring that \(u_{n-1}\) and \(u_n\) live in
    different coordinates.

    \item \textbf{Joint mean-field with Gaussian projected-prior transfer.}
    This is the efficient filtering-style approximation.

    \item \textbf{Joint mean-field with SSGP-style old-likelihood-ratio transfer.}
    This is the main proposed streaming variational method.

    \item \textbf{Structured joint posterior with SSGP-style transfer.}
    This tests whether retaining part of the \(\beta\)-\(u\) posterior dependence improves
    calibration.
\end{enumerate}

\subsection{Ablations}

Important ablations include:
\begin{enumerate}[label=(\arabic*)]
    \item number of temporal HiPPO basis functions \(M_t\);
    \item number of spatial inducing points \(M_s\);
    \item number of RFF frequencies \(R\);
    \item fixed temporal basis versus changing HiPPO basis;
    \item no transfer versus projected-prior transfer versus old-likelihood-ratio transfer;
    \item exact old-likelihood transfer versus structured approximation to
    \(\Lambda_{o\rightarrow n}\);
    \item maintaining \(R_o\) as a dense matrix versus maintaining \(B_o\otimes G\);
    \item mean-field posterior versus full or structured joint posterior;
    \item static \(\beta\) versus drifting \(\beta\) with \(Q_\beta\);
    \item PSD projection, damping, or jitter for \(R_o=S_o^{-1}-K_{oo}^{-1}\) when a
    dense approximation is used.
\end{enumerate}

\subsection{Metrics}

The evaluation should not rely only on RMSE. Report:
\begin{enumerate}[label=(\arabic*)]
    \item RMSE;
    \item MAE;
    \item test negative log likelihood;
    \item predictive interval coverage, for example \(50\%\), \(90\%\), and \(95\%\);
    \item runtime per online block;
    \item peak memory;
    \item RMSE/NLL under matched memory budgets;
    \item calibration error for Gaussian predictive intervals;
    \item error on continual evaluation windows, not only future-only forecasting.
\end{enumerate}

The most important comparison is whether SSGP-style joint training improves continual
evaluation and uncertainty calibration relative to residual-based two-stage learning and
relative to the simpler projected-prior transfer.

\section{Summary}
\label{sec:summary}

The proposed model follows the logical structure
\begin{equation}
\begin{aligned}
    &\text{two-stage learning has stale residuals} \\
    \Rightarrow\;&
    \text{joint training fixes the observation target} \\
    \Rightarrow\;&
    \text{SVGP ELBO gives the variational framework} \\
    \Rightarrow\;&
    \text{HiPPO changing basis requires old-information transfer} \\
    \Rightarrow\;&
    \text{SSGP-style old-likelihood-ratio transfer gives a principled streaming update} \\
    \Rightarrow\;&
    \text{Kronecker structure makes the method scalable.}
\end{aligned}
\end{equation}

The final online observation model is
\begin{equation}
    y_n
    =
    \Phi_n\beta
    +
    A_nu_n
    +
    \epsilon_n,
    \qquad
    \epsilon_n\sim\mathcal{N}(0,\sigma^2I).
\end{equation}
The key distinction from residual-based learning is that the likelihood is always applied
to \(y_n\), not to a residual target that changes whenever \(\beta\) is updated:
\begin{equation}
    p(y_n\mid \beta,u_n)
    =
    \mathcal{N}
    \left(
        y_n
        \mid
        \Phi_n\beta+A_nu_n,
        \sigma^2I
    \right).
\end{equation}

For the linear coefficient, the streaming prior is
\begin{equation}
    p_n(\beta)=q_{n-1}(\beta),
\end{equation}
or, with slow drift,
\begin{equation}
    p_n(\beta)
    =
    \mathcal{N}
    \left(
        m_{\beta,n-1},
        S_{\beta,n-1}+Q_\beta
    \right).
\end{equation}

For the GP inducing variables under a changing HiPPO temporal basis, the main proposed
update follows the streaming sparse GP principle. Instead of directly reusing
\(q_{n-1}(u_o)\), we propagate the old likelihood ratio
\begin{equation}
    \frac{q_{n-1}(u_o)}{p(u_o)}.
\end{equation}
For Gaussian \(q_{n-1}(u_o)=\mathcal{N}(m_o,S_o)\) and
\(p(u_o)=\mathcal{N}(0,K_{oo})\), define
\begin{equation}
    R_o=S_o^{-1}-K_{oo}^{-1},
    \qquad
    r_o=S_o^{-1}m_o.
\end{equation}
Using
\begin{equation}
    p(u_o\mid u_n)
    =
    \mathcal{N}
    \left(
        L_{on}u_n,
        \Sigma_{o\mid n}
    \right),
    \qquad
    L_{on}=K_{on}K_{nn}^{-1},
\end{equation}
the old likelihood information is transferred to the new HiPPO inducing coordinates
through
\begin{equation}
    \Lambda_{o\rightarrow n}
    =
    L_{on}^\top R_oL_{on},
    \qquad
    h_{o\rightarrow n}
    =
    L_{on}^\top r_o.
\end{equation}
The main GP coordinate update is therefore
\begin{equation}
    S_{u,n}^{-1}
    =
    K_{nn}^{-1}
    +
    \Lambda_{o\rightarrow n}
    +
    \frac{1}{\sigma^2}A_n^\top A_n,
\end{equation}
and
\begin{equation}
    S_{u,n}^{-1}m_{u,n}
    =
    h_{o\rightarrow n}
    +
    \frac{1}{\sigma^2}
    A_n^\top
    \left(
        y_n-\Phi_nm_{\beta,n}
    \right).
\end{equation}

Crucially, this SSGP-style transfer can preserve the Sylvester structure under the
Kronecker-separable HiPPO-SVGP assumptions. Since the changing basis acts only in the
temporal dimension and the spatial inducing set is fixed,
\begin{equation}
    L_{on}=L_{on}^{(t)}\otimes I_s.
\end{equation}
If the old likelihood precision is maintained as
\begin{equation}
    R_o=B_o\otimes G,
    \qquad
    G=C^\top C,
\end{equation}
then
\begin{equation}
    \Lambda_{o\rightarrow n}
    =
    \left[
        \left(L_{on}^{(t)}\right)^\top
        B_o
        L_{on}^{(t)}
    \right]\otimes G.
\end{equation}
Thus the updated precision keeps the form
\begin{equation}
    \Lambda_{u,n}
    =
    \left(K_{nn}^{(t)}\right)^{-1}
    \otimes K_s^{-1}
    +
    B_n\otimes G,
\end{equation}
with
\begin{equation}
    B_n
    =
    \left(L_{on}^{(t)}\right)^\top
    B_o
    L_{on}^{(t)}
    +
    \frac{1}{\sigma^2}T_n^\top T_n.
\end{equation}
Therefore, the Sylvester solver remains applicable:
\begin{equation}
    K_s^{-1}Z
    \left(K_{nn}^{(t)}\right)^{-1}
    +
    GZB_n
    =
    Q.
\end{equation}

The Gaussian projected-prior update
\begin{equation}
    p_n^{\mathrm{proj}}(u_n)
    =
    \int
    p(u_n\mid u_o)q_{n-1}(u_o)\,du_o
\end{equation}
is retained as an efficient approximation and ablation baseline. The main theoretical
contribution, however, is the SSGP-style old-likelihood-ratio transfer together with the
Kronecker-preserving temporal transfer statistic. This combination gives a principled
streaming variational method for joint linear-mean HiPPO-SVGP while preserving the
scalable Sylvester computation required for large spatio-temporal inducing states.


\appendix

\section{Proofs for Kronecker-Preserving Changing-Basis Transfer}
\label{app:kron_preserving_transfer_proofs}

This appendix gives the detailed proof of the two structural claims used
in the main text:
\begin{equation}
    L_{on}=L_{on}^{(t)}\otimes I_s,
    \label{eq:app_Lon_kron_claim}
\end{equation}
and
\begin{equation}
    R_o=B_o\otimes G.
    \label{eq:app_Ro_kron_claim}
\end{equation}
The first identity is a consequence of the separable spatio-temporal
kernel and the fact that the spatial inducing set is fixed. The second
identity is not true for an arbitrary dense posterior covariance; it is an
algorithmic invariant obtained by maintaining the old likelihood natural
precision in a Kronecker-structured form.

\subsection{\texorpdfstring{Proof that \(L_{on}=L_{on}^{(t)}\otimes I_s\)}{Proof that Lon = Lon(t) tensor Is}}
\label{app:proof_Lon_temporal_only}

Assume the spatio-temporal kernel is separable:
\begin{equation}
    k\big((t,s),(t',s')\big)
    =
    k_t(t,t')k_s(s,s').
    \label{eq:app_sep_kernel}
\end{equation}
Let the mixed HiPPO interdomain inducing variables be
\begin{equation}
    u_{\ell,j}(T)
    =
    \int_0^T
    g_\ell^{(T)}(x)f(x,z_j^{(s)})\,dx,
    \label{eq:app_mixed_inducing}
\end{equation}
where \(\ell\) indexes the temporal HiPPO basis and \(j\) indexes the
fixed spatial inducing location \(z_j^{(s)}\).

Let
\begin{equation}
    u_o := u(T_o),
    \qquad
    u_n := u(T_n),
\end{equation}
where \(T_o=T_{n-1}\) is the old temporal horizon and \(T_n\) is the new
temporal horizon. The temporal HiPPO basis changes from
\(g_\ell^{(T_o)}\) to \(g_\ell^{(T_n)}\), but the spatial inducing
locations \(Z_s=\{z_1^{(s)},\ldots,z_{M_s}^{(s)}\}\) are fixed.

For two inducing entries,
\begin{equation}
    u_{\ell,j}^{old}
    =
    \int_0^{T_o}
    g_\ell^{(T_o)}(x)f(x,z_j^{(s)})\,dx,
\end{equation}
and
\begin{equation}
    u_{m,j'}^{new}
    =
    \int_0^{T_n}
    g_m^{(T_n)}(x')f(x',z_{j'}^{(s)})\,dx',
\end{equation}
their prior covariance is
\begin{align}
    \cov
    \left(
        u_{\ell,j}^{old},
        u_{m,j'}^{new}
    \right)
    &=
    \int_0^{T_o}
    \int_0^{T_n}
    g_\ell^{(T_o)}(x)
    g_m^{(T_n)}(x')
    k\big((x,z_j^{(s)}),(x',z_{j'}^{(s)})\big)
    \,dx'\,dx                                                  \\
    &=
    \int_0^{T_o}
    \int_0^{T_n}
    g_\ell^{(T_o)}(x)
    g_m^{(T_n)}(x')
    k_t(x,x')
    k_s(z_j^{(s)},z_{j'}^{(s)})
    \,dx'\,dx                                                   \\
    &=
    \left[
    \int_0^{T_o}
    \int_0^{T_n}
    g_\ell^{(T_o)}(x)
    g_m^{(T_n)}(x')
    k_t(x,x')
    \,dx'\,dx
    \right]
    k_s(z_j^{(s)},z_{j'}^{(s)}).
    \label{eq:app_old_new_cov_entry}
\end{align}
Thus the old-new covariance factorizes into a temporal covariance factor
and a spatial covariance factor. With the vectorization convention used in
the main text,
\begin{equation}
    K_{on}
    =
    K_{on}^{(t)}\otimes K_s,
    \qquad
    K_s:=K_{ZZ}^{(s)}.
    \label{eq:app_Kon_kron}
\end{equation}
Similarly,
\begin{equation}
    K_{nn}
    =
    K_{nn}^{(t)}\otimes K_s,
    \qquad
    K_{oo}
    =
    K_{oo}^{(t)}\otimes K_s.
    \label{eq:app_Knn_Koo_kron}
\end{equation}

The conditional mean operator in
\begin{equation}
    p(u_o\mid u_n)
    =
    \mathcal{N}
    \left(
        L_{on}u_n,
        \Sigma_{o\mid n}
    \right)
\end{equation}
is
\begin{equation}
    L_{on}
    =
    K_{on}K_{nn}^{-1}.
\end{equation}
Using \eqref{eq:app_Kon_kron}--\eqref{eq:app_Knn_Koo_kron} and the
Kronecker identities
\begin{equation}
    (A\otimes B)^{-1}=A^{-1}\otimes B^{-1},
    \qquad
    (A\otimes B)(C\otimes D)=(AC)\otimes(BD),
\end{equation}
we obtain
\begin{align}
    L_{on}
    &=
    \left(
        K_{on}^{(t)}\otimes K_s
    \right)
    \left[
        \left(K_{nn}^{(t)}\right)^{-1}
        \otimes
        K_s^{-1}
    \right]                                                     \\
    &=
    K_{on}^{(t)}
    \left(K_{nn}^{(t)}\right)^{-1}
    \otimes
    K_sK_s^{-1}                                                  \\
    &=
    L_{on}^{(t)}\otimes I_s,
    \label{eq:app_Lon_kron_proof}
\end{align}
where
\begin{equation}
    L_{on}^{(t)}
    :=
    K_{on}^{(t)}
    \left(K_{nn}^{(t)}\right)^{-1}.
\end{equation}
Therefore, under a separable kernel, fixed spatial inducing locations, and
temporal-only changing basis, the old-to-new conditional operator acts
only on the temporal inducing coordinates and leaves the spatial inducing
coordinates unchanged.

\paragraph{Vectorization convention.}
The identity \eqref{eq:app_Lon_kron_proof} assumes the same vectorization
ordering as the Kronecker convention \(K_{uu}=K_{uu}^{(t)}\otimes K_s\).
If an implementation uses the opposite vectorization order, the equivalent
identity becomes \(L_{on}=I_s\otimes L_{on}^{(t)}\). The mathematics is
unchanged, but the code must be consistent with the chosen ordering.

\subsection{\texorpdfstring{Proof and Status of \(R_o=B_o\otimes G\)}{Proof and Status of Ro = Bo tensor G}}
\label{app:proof_Ro_kron_natural_stat}

The identity
\begin{equation}
    R_o=B_o\otimes G
\end{equation}
should be interpreted carefully. It is not a property of an arbitrary
dense posterior covariance. It is a structured natural-parameter
representation of the accumulated old likelihood information.

For a Gaussian observation model, the likelihood contribution of a block
\(j\) to the inducing-variable precision is
\begin{equation}
    \Lambda_{\mathrm{like},j}
    =
    \frac{1}{\sigma^2}A_j^\top A_j.
    \label{eq:app_likelihood_precision}
\end{equation}
Under the Kronecker HiPPO-SVGP representation,
\begin{equation}
    A_j=T_j\otimes C,
    \label{eq:app_Aj_kron}
\end{equation}
where \(T_j\) is the temporal projection matrix for block \(j\), and
\begin{equation}
    C=K_{XZ}^{(s)}\left(K_{ZZ}^{(s)}\right)^{-1}
\end{equation}
is the fixed spatial projection matrix. Therefore,
\begin{align}
    \Lambda_{\mathrm{like},j}
    &=
    \frac{1}{\sigma^2}
    (T_j\otimes C)^\top(T_j\otimes C)                         \\
    &=
    \frac{1}{\sigma^2}
    (T_j^\top T_j)\otimes(C^\top C).
    \label{eq:app_likelihood_precision_kron}
\end{align}
Define
\begin{equation}
    G:=C^\top C.
\end{equation}
Then each block contributes
\begin{equation}
    \Lambda_{\mathrm{like},j}
    =
    \frac{1}{\sigma^2}
    (T_j^\top T_j)\otimes G.
\end{equation}

The old likelihood natural precision before processing block \(n\) is the
sum of historical likelihood precision contributions, expressed in the old
inducing coordinates:
\begin{equation}
    R_o
    :=
    \sum_{j<n}
    \Lambda_{\mathrm{like},j}.
\end{equation}
If the spatial projection matrix \(C\) is fixed across the online blocks,
then the spatial factor \(G=C^\top C\) is common to all terms. Hence
\begin{align}
    R_o
    &=
    \sum_{j<n}
    \frac{1}{\sigma^2}
    (T_j^\top T_j)\otimes G                                  \\
    &=
    \left[
        \sum_{j<n}
        \frac{1}{\sigma^2}
        T_j^\top T_j
    \right]\otimes G.
\end{align}
Define the temporal old-likelihood precision statistic
\begin{equation}
    B_o
    :=
    \sum_{j<n}
    \frac{1}{\sigma^2}
    T_j^\top T_j.
    \label{eq:app_Bo_definition}
\end{equation}
Then
\begin{equation}
    R_o=B_o\otimes G.
    \label{eq:app_Ro_Bo_G}
\end{equation}

Thus \(R_o=B_o\otimes G\) is valid when the algorithm maintains the old
likelihood precision directly in natural-parameter form. In exact
Gaussian natural-parameter inference, this agrees with
\begin{equation}
    R_o=S_o^{-1}-K_{oo}^{-1},
\end{equation}
provided \(S_o^{-1}\) denotes the corresponding inducing-variable
posterior precision in the same structured representation. However, if
one stores an unrestricted dense posterior covariance, a marginal
covariance from a full joint posterior over \((\beta,u_o)\), or an
arbitrary structured approximation, then the dense quantity
\(S_o^{-1}-K_{oo}^{-1}\) need not be a single Kronecker product.

\subsection{Kronecker Preservation of the Old-Likelihood Transfer}
\label{app:proof_Lambda_transfer_kron}

Combining the previous two results proves the Kronecker preservation of
the SSGP-style old-likelihood-ratio transfer.

The transferred old likelihood precision is
\begin{equation}
    \Lambda_{o\rightarrow n}
    =
    L_{on}^\top R_oL_{on}.
\end{equation}
Using
\begin{equation}
    L_{on}=L_{on}^{(t)}\otimes I_s,
    \qquad
    R_o=B_o\otimes G,
\end{equation}
we have
\begin{align}
    \Lambda_{o\rightarrow n}
    &=
    \left[
        \left(L_{on}^{(t)}\right)^\top\otimes I_s
    \right]
    (B_o\otimes G)
    \left[
        L_{on}^{(t)}\otimes I_s
    \right]                                                     \\
    &=
    \left[
        \left(L_{on}^{(t)}\right)^\top
        B_o
        L_{on}^{(t)}
    \right]
    \otimes G.
\end{align}
Define
\begin{equation}
    B_{o\rightarrow n}
    :=
    \left(L_{on}^{(t)}\right)^\top
    B_o
    L_{on}^{(t)}.
\end{equation}
Then
\begin{equation}
    \Lambda_{o\rightarrow n}
    =
    B_{o\rightarrow n}\otimes G.
    \label{eq:app_Lambda_transfer_kron}
\end{equation}
The new block likelihood contribution is
\begin{equation}
    \frac{1}{\sigma^2}A_n^\top A_n
    =
    \frac{1}{\sigma^2}
    (T_n^\top T_n)\otimes G.
\end{equation}
Therefore the updated temporal likelihood statistic is
\begin{equation}
    B_n
    =
    B_{o\rightarrow n}
    +
    \frac{1}{\sigma^2}T_n^\top T_n.
    \label{eq:app_Bn_after_transfer}
\end{equation}
The resulting inducing precision has the Sylvester-compatible form
\begin{equation}
    \Lambda_{u,n}
    =
    \left(K_{nn}^{(t)}\right)^{-1}\otimes K_s^{-1}
    +
    B_n\otimes G.
    \label{eq:app_Lambda_un_sylvester_form}
\end{equation}

Consequently, for \(q_*=\vecop(Q)\) and
\(z=\Lambda_{u,n}^{-1}q_*=\vecop(Z)\), the system
\begin{equation}
    \Lambda_{u,n}z=q_*
\end{equation}
is equivalent to the Sylvester-type equation
\begin{equation}
    K_s^{-1}Z
    \left(K_{nn}^{(t)}\right)^{-1}
    +
    GZB_n
    =
    Q,
    \label{eq:app_sylvester_after_transfer}
\end{equation}
assuming the vectorization convention used in the main text and symmetric
\(B_n\). Thus the SSGP-style old-likelihood-ratio transfer can preserve
the Sylvester solver under the temporal-only changing-basis and fixed
spatial-geometry assumptions.

\subsection{\texorpdfstring{When \(R_o=B_o\otimes G\) Can Fail}{When Ro = Bo tensor G Can Fail}}
\label{app:failure_cases_Ro_kron}

The single-Kronecker representation \(R_o=B_o\otimes G\) is not
unconditional. It may fail in the following cases.

\paragraph{Changing spatial observation pattern.}
If the spatial observation pattern or spatial mini-batch changes across
blocks, then the spatial projection matrix may become block-dependent:
\begin{equation}
    C_j\neq C.
\end{equation}
Consequently,
\begin{equation}
    G_j=C_j^\top C_j
\end{equation}
also changes. The old likelihood precision becomes
\begin{equation}
    R_o
    =
    \sum_{j<n}
    B_j\otimes G_j,
\end{equation}
rather than a single Kronecker product \(B_o\otimes G\). This is still a
sum-of-Kronecker structure, but the simple closed-form Sylvester equation
in \eqref{eq:app_sylvester_after_transfer} no longer applies directly.
One would need a generalized Sylvester solver or an iterative method using
Kronecker matrix-vector products.

\paragraph{Changing spatial inducing locations.}
If the spatial inducing locations \(Z_s\) are updated online, then
\begin{equation}
    K_s=K_{ZZ}^{(s)},
    \qquad
    C=K_{XZ}^{(s)}\left(K_{ZZ}^{(s)}\right)^{-1},
    \qquad
    G=C^\top C
\end{equation}
all change over time. In that case, the old-to-new transfer operator need
not have the form \(L_{on}^{(t)}\otimes I_s\), and the likelihood
precision is generally not a single \(B_o\otimes G\) term.

\paragraph{Recovering \(R_o\) from a dense posterior covariance.}
If one stores a full joint posterior over \((\beta,u_o)\) and then uses
only the marginal covariance of \(u_o\), or if one stores an unrestricted
dense posterior covariance \(S_o\) and directly computes
\begin{equation}
    R_o=S_o^{-1}-K_{oo}^{-1},
\end{equation}
then \(R_o\) is generally not a single Kronecker product unless the
posterior precision itself preserves the required structure. For the
Sylvester-preserving implementation, the old likelihood natural precision
should be maintained directly as \(B_o\otimes G\), rather than recovered
from an arbitrary dense covariance.

\paragraph{Non-Gaussian likelihoods or non-conjugate approximations.}
For Gaussian likelihoods, the likelihood precision contribution is
\begin{equation}
    \frac{1}{\sigma^2}A_j^\top A_j.
\end{equation}
With non-Gaussian likelihoods, approximate inference methods such as
Laplace, expectation propagation, or stochastic variational inference may
produce likelihood-site precisions that are not simply proportional to
\(A_j^\top A_j\), or that vary by observation. The resulting precision
may fail to share a common spatial factor \(G\).

\paragraph{Practical implication.}
To retain the simple Sylvester solver, the scalable version of the method
should explicitly assume fixed spatial inducing geometry and should
maintain the old likelihood precision as a Kronecker natural statistic:
\begin{equation}
    R_o=B_o\otimes G.
\end{equation}
When these assumptions are violated, the model can still be implemented,
but the solver must be replaced by a more general structured linear
algebra routine, such as a sum-of-Kronecker conjugate-gradient method or a
generalized Sylvester solver.

\end{document}
